\documentclass[11pt,oneside]{report}
% ======================================================================
%  THE TRIADIC MEASUREMENT SUBSTRATE --- COMPLETE CORPUS (Volumes 1 & 2)
%  Aaron Switzer, 2026
%  Single-file build. Compile: pdflatex x3 (third pass settles the TOC).
% ======================================================================
% ======================================================================
%  TMS Volume 1 -- shared preamble
%  Compiles with pdfLaTeX both locally and on Overleaf (no exotic fonts).
% ======================================================================
\usepackage[T1]{fontenc}
\usepackage[utf8]{inputenc}
\usepackage{amsmath}
\usepackage{amssymb}
\usepackage{mathptmx}            % Times text + math (widely available)
\usepackage[scaled=0.90]{helvet} % sans for headings
\usepackage{courier}

\usepackage[letterpaper,margin=1in]{geometry}
\usepackage{amsthm}
\usepackage{mathtools}
\usepackage{bm}
\usepackage{microtype}
\usepackage{enumitem}
\usepackage{booktabs}
\usepackage{array}
\usepackage{longtable}
\usepackage{caption}
\usepackage{xcolor}
\usepackage{titlesec}
\usepackage{fancyhdr}
\usepackage[hidelinks]{hyperref}
\usepackage{tocbibind}

% ---- spacing -----------------------------------------------------------
\setlength{\parindent}{1.2em}
\setlength{\parskip}{0.35em}
\linespread{1.04}
\setlist{leftmargin=1.6em,topsep=0.25em,itemsep=0.15em,parsep=0pt}

% ---- theorem environments ---------------------------------------------
\newtheorem{theorem}{Theorem}[section]
\newtheorem{lemma}[theorem]{Lemma}
\newtheorem{corollary}[theorem]{Corollary}
\newtheorem{proposition}[theorem]{Proposition}

\theoremstyle{definition}
\newtheorem{definition}[theorem]{Definition}
\newtheorem{axiom}{Axiom}
\newtheorem{claim}[theorem]{Claim}
\newtheorem{principle}[theorem]{Principle}
\newtheorem{conjecture}[theorem]{Conjecture}
\newtheorem{prediction}[theorem]{Prediction}
\newtheorem{property}[theorem]{Property}
\newtheorem{postulate}[theorem]{Postulate}

\theoremstyle{remark}
\newtheorem{remark}[theorem]{Remark}
\newtheorem*{remark*}{Remark}

% ---- section heading look (unnumbered; author supplies the numbers) ----
\titleformat{\section}[hang]{\normalfont\large\bfseries\sffamily}{}{0pt}{}
\titleformat{\subsection}[hang]{\normalfont\normalsize\bfseries\sffamily}{}{0pt}{}
\titleformat{\subsubsection}[hang]{\normalfont\normalsize\bfseries\itshape}{}{0pt}{}
\titlespacing*{\section}{0pt}{1.4em}{0.5em}
\titlespacing*{\subsection}{0pt}{1.0em}{0.35em}
\titlespacing*{\subsubsection}{0pt}{0.8em}{0.3em}

% chapter (= each paper) heading
\titleformat{\chapter}[display]
  {\normalfont\sffamily\bfseries}
  {}{0pt}{\Huge}
\titlespacing*{\chapter}{0pt}{10pt}{24pt}

% ---- page-break quality: avoid stranded headings and orphaned/widowed lines ----
% (applies to both volumes via the shared preamble)
\widowpenalty=10000        % no single line at the top of a page
\clubpenalty=10000         % no single line at the bottom of a page
\displaywidowpenalty=10000 % same, around displayed equations
\brokenpenalty=4000        % discourage page breaks right after a hyphenated line
% titlesec: forbid a page break immediately after a heading (heading stranded at page foot)
\usepackage{needspace}
\usepackage{etoolbox}
\pretocmd{\section}{\Needspace*{4\baselineskip}}{}{}
\pretocmd{\subsection}{\Needspace*{3\baselineskip}}{}{}

% ---- running heads -----------------------------------------------------
\pagestyle{fancy}
\fancyhf{}
\fancyhead[L]{\small\itshape The Triadic Measurement Substrate}
\fancyhead[R]{\small\itshape\nouppercase{\rightmark}}
\fancyfoot[C]{\small\thepage}
\renewcommand{\headrulewidth}{0.4pt}
\fancypagestyle{plain}{\fancyhf{}\fancyfoot[C]{\small\thepage}\renewcommand{\headrulewidth}{0pt}}

% ---- abstract / keywords / references list ----------------------------
\newenvironment{paperabstract}
  {\begin{center}\begin{minipage}{0.88\textwidth}\small
   \noindent\textbf{Abstract.}\itshape\space}
  {\end{minipage}\end{center}\vspace{0.4em}}

\newcommand{\keywords}[1]{\noindent{\small\textbf{Keywords:} #1}\vspace{0.4em}}

% per-paper APA reference list (hanging indent), kept as authored
\newenvironment{reflist}
  {\par\small\setlength{\parindent}{0pt}%
   \setlength{\leftskip}{1.6em}\setlength{\parindent}{-1.6em}%
   \setlength{\parskip}{0.4em}%
   \raggedright\hyphenpenalty=50\exhyphenpenalty=50\relax
   \sloppy}
  {\par}

% convenience
\providecommand{\tightlist}{\setlength{\itemsep}{0pt}\setlength{\parskip}{0pt}}
\providecommand{\TMS}{\textsc{tms}}
\hyphenation{con-fir-ma-tion sub-strate}
\begin{document}

\thispagestyle{empty}
\begin{titlepage}
\centering
\vspace*{2cm}
{\sffamily\bfseries\Large The Triadic Measurement Substrate\par}
\vspace{0.6cm}
{\large $\bar{B} = 3\bar{\alpha}$\par}
\vspace{1cm}
{\sffamily\large Volume 2\par}
\vspace{1.4cm}
{\Large \bfseries Paper 4: Time: The Now-Surface Deepened --- Confirmed Past, Unwritten Future, and the Decay of the Wake\par}
\vfill
{\large Aaron Switzer\par}
\vspace{0.4cm}
{\large 2026\par}
\vspace{1.2cm}
{\small\itshape Excerpted from the complete two-volume corpus, \emph{The Triadic Measurement Substrate}. Full corpus and reproducibility package: \url{https://doi.org/10.5281/zenodo.22035422}\par}
\end{titlepage}
\cleardoublepage
\pagenumbering{arabic}

% ---------------- Volume 2 / P4_Time ----------------
\chapter*{Paper 4}
\markboth{Paper 4}{Paper 4}
\addcontentsline{toc}{chapter}{Paper 4: Time: The Now-Surface Deepened --- Confirmed Past, Unwritten Future, and the Decay of the Wake}
\begingroup\centering
{\large\itshape Time: The Now-Surface Deepened --- Confirmed Past, Unwritten Future, and the Decay of the Wake\par}\vspace{0.8em}
{\normalsize Aaron Switzer\par}
{\small June 2026\par}
\endgroup\vspace{1.0em}

\noindent{\small\textbf{Keywords:} Now-Surface, Wake, Confirmed Past, Unwritten Future, Reinterpretation Operator, Decay, Forgetting, Re-read Clock, Excitable Manifold, Rate-Distortion, Triadic Confirmation}\par\vspace{0.6em}

\section*{Status of Results}
\addcontentsline{toc}{section}{Status of Results}

This paper continues the labeling conventions of Volume 1 and of the
preceding papers of this volume. Every claim in the body falls into
exactly one category, marked where it appears. No claim is upgraded
beyond what its derivation earns.

\textbf{Derived.} Results that follow from the five axioms, the Principle
of Generative Economy, and \ensuremath{U_{\mathrm{sh}}} by explicit
construction or proof, or from results already derived in Volume~1 and
the earlier papers of this volume by explicit inheritance. The
asymmetry between a confirmed past and an unwritten future as the two
sides of the now-surface (\S\,I); the reconstructability of an interior
confirmation's value from the conserved polarization it deposits into its
forward cone, increasing with depth into the wake, with detection-only at
the frontier (\S\,I); the fixed-fraction form of re-read decay, forced by
the access constraint on the reinterpretation operator (\S\,IV,
Appendix~A); the monotone forgetting curve with a count-conserved floor
(\S\,IV); and the structural flattening of the cosmic-residue balance
(\S\,IV) are derived in this sense.

\textbf{Derived-conditional.} Results that follow given a primitive
established elsewhere whose own grounding is stated. The per-fold shed
value \ensuremath{s = 2} is inherited from Paper~3, where it is forced
given binary-agent triads; this paper uses the value and does not
re-establish it.

\textbf{Structural correspondence.} Results in which a TMS construction
reproduces the skeleton of a known framework or quantifies a recognized
distinction, never used as support for a TMS claim, only as a landing
point after one. The reading of re-read decay as lossy recompression in
the rate-distortion sense (\S\,IV, \S\,V); and the correspondence of the
poised pre-confirmation configuration to an excitable system held near
threshold (\S\,III, \S\,V) are structural correspondences in this sense.

\textbf{Predictions / Conjectures.} Falsifiable claims that follow from
the framework but whose empirical tests are not established here. That a
memory degrades by the same law whether overwritten from the world or
re-read from within --- that memory decays through use, not only through
neglect (\S\,III) --- is a prediction in this sense; its match to data is
not claimed.

\textbf{Permanent open horizon.} A question the framework is structurally
unable to close, marked so that no later claim assumes it closed. What
fills an unwritten interstice when it is finally confirmed --- whether the
content of the future has any determination prior to its folding --- is
not settled by anything in this paper and is left open (\S\,I, \S\,II).
This silence is load-bearing.

\textbf{Deferred.} Quantities the framework defines but does not evaluate
here. The value of the cosmic residue-to-ordinary ratio depends on two
rates this paper does not derive and is deferred (\S\,IV). The
identification of the anchored residue with the dark-matter mode, and the
mass--energy reading that binds it to dark energy, belong to the material
side and are treated there, not here (\S\,IV).

\section*{Notation}

Symbols are inherited from Volume~1 and the earlier papers of this
volume. \ensuremath{\hat{R}} is the reinterpretation operator (Volume~1,
Definition~0.0.67); the now-surface is the active codimension-one
confirmation frontier and the wake is the settled three-dimensional
record behind it (Paper~1, \S\,1); \ensuremath{U_{\mathrm{sh}}} is the
maximum-entropy substrate field; \ensuremath{s} is the per-fold
distinction shed of Paper~3, with \ensuremath{s = 2} for binary-agent
triads; \ensuremath{n} counts folds applied to a given structure;
\ensuremath{r = 2^{-s}} is the per-fold survival ratio introduced in
\S\,IV.

\section*{I. The Now-Surface Deepened}
\addcontentsline{toc}{section}{I. The Now-Surface Deepened}

Paper~1 introduced the now-surface as the active frontier where agents
are presently locking bits, and the wake as the settled record that
accumulates behind it. That construction was developed entirely on the
material side; it carried gravity, the black hole, and dynamic dark
energy, and assumed no property of observers. This paper takes the same
object and asks what it determines about time itself. The question is
single: where, on the substrate, is the asymmetry between past and
future, and what governs the relation of one now to the next.

The answer is already present in the three-state space of Volume~1.
Paper~1 \S\,1.4 fixes three ontologically distinct states for an
interstice: latent, drawn from \ensuremath{U_{\mathrm{sh}}} and neither
zero nor one but the incompressible superposition of both;
confirmed-zero; and confirmed-one. Confirmation resolves the latent state
into one of the two confirmed values, and by Axiom~4 the result cannot be
overwritten. These three states, read against the now-surface, are the
three regions of time.

The confirmed wake is the past. It is the locked record, sealed by
Axiom~4: the system has only one available direction, and each flop adds
a layer without disturbing prior layers. The now-surface is the present:
the codimension-one frontier where confirmation is currently occurring,
the only place a fold can fire. The future is the latent population the
present surface has not yet folded. The identification is relative to a
surface and not global: latency is a property an interstice has at large,
but it reads as future only with respect to a given now-surface, for
which it is the not-yet-confirmed region that surface can advance into.
What is future to one frontier is unrelated to another that will never
reach it. So read, the future is unwritten in the strict sense: not
merely unknown but undetermined, the superposition not yet resolved into
a confirmed value.

This locates the asymmetry exactly. The past and the future are not two
symmetric directions along a line; they are two different states of an
interstice relative to a single surface. The past is what has been
confirmed and sealed; the future is what remains latent; the now-surface
is the boundary between them, and it advances under \ensuremath{\hat{R}}
every flop (Paper~1, \S\,7). The arrow of time is not added to this
picture. It is Axiom~4 read as a statement about which side of the
now-surface a fold can act on. A fold acts in exactly two ways. It can
confirm a latent interstice, resolving the superposition into a confirmed
value and moving the interstice from the future side of the surface to
the past side. It can also re-read a region of wake already confirmed.
What it cannot do is return a confirmed interstice to the latent state:
that is precisely the overwrite Axiom~4 forbids. Confirmation therefore
runs one way across the surface, and the one-way-ness is not a separate
postulate but the causal anchoring of Volume~1 read at the surface.

Two consequences follow immediately, and they organize the rest of the
paper. First, only the past decays. A latent interstice carries no
confirmed distinction to lose; there is nothing in the future to shed,
because nothing there has been written. Degradation is therefore a
strictly past-facing phenomenon: it acts on the wake, never on the latent
population. The temporal arrow and the direction of decay are the same
fact --- one can re-read the wake and lose grain, but one cannot re-read
the unwritten, because there is nothing yet to read. Decay is developed in
\S\,IV.

The asymmetry is directional, not an asymmetry of determination. The
substrate geometry runs both ways across the surface, and the difference
between past and future is the difference between running it backward into
the wake and running it forward into the latent region. Run backward, the
geometry reconstructs, and it does so for a reason internal to the
substrate. Each confirmation transports its polarization outward through
the \ensuremath{1\to 4} multiplier: a confirmed value emits its signals
into the downstream synthetic bits, which carry that value's polarization
forward, and by Axiom~3 those signals are conserved and cannot be lost
(Paper~1, \S\,4.2--4.3). A confirmation in the interior of the wake has
therefore already deposited its value, redundantly, across the layers it
produced. If a single agent's contribution is later disrupted, the local
unit only registers that a confirmation was present --- the boundary opens
and a parity syndrome appears --- but the original value remains
recoverable from the conserved polarization the confirmation laid into its
forward cone. The past is rigid because it is over-determined by its own
downstream record, and the over-determination grows with depth: the deeper
a confirmation lies in the wake, the more forward layers redundantly
encode it, so older structure is more strongly fixed than recent
structure. This redundancy structure corresponds to the \ensuremath{[[4,2,2]]}
error-detection code, whose single-error detection sits on the local unit
and whose reconstruction is carried by the forward cone (Paper~1, \S\,5);
the correspondence is noted as a landing, not relied on as the argument.

The reconstruction has one limit, and it falls exactly at the present. A
confirmation at the now-surface frontier has not yet produced a downstream
cone, so a disruption there can only be detected and re-confirmed against
the current polarization landscape, which biases the redrawn value without
fixing it. Rigidity of value is a property of the wake interior; at the
frontier only the fact of confirmation is preserved, not its value. This
is consistent with the asymmetry as a whole: the past is rigid in
proportion to how much wake lies downstream of it, and the present, having
none, is the one place the record is not yet fixed.

Run forward, the same geometry is blocked, and by two obstructions of
different kinds. The first is the latent draw itself: the next
confirmation resolves against a fresh contribution from
\ensuremath{U_{\mathrm{sh}}}, which is maximum-entropy and has no
attractors (Paper~1, Axiom~0; \S\,7.2), so the forward step is not
determined by present structure however completely that structure is
known. This obstruction is ontological and holds for any observer. The
second is positional: a bounded subsystem cannot measure across its own
closure boundary, and the exterior it cannot read co-determines the folds
it will undergo, so from inside the subsystem the forward step is
unreadable even where it is determined. This obstruction is relative to
the observer's position with respect to the boundary; an observer outside
the system would not encounter it. The forward direction is therefore
unshortcuttable from within, which is the time-like face of the
record's computational irreducibility (Paper~1, \S\,7.2a). The past is
reconstructable and the future is not, from one geometry, because
reconstruction runs with the closed constraints and the conserved
downstream record, while prediction runs against an undetermined draw and
an unreadable exterior.

Second, the future is not modeled by reaching toward it. Nothing on the
present surface extends into the latent region or carries information
about which value an interstice will take. What a system can do is fold
on its own wake --- on structure already written --- to build a present
configuration, and that configuration can bias which latent interstices
the surface confirms next without containing any of their content. The
configuration is past-built and present-held; it points at the future
only in that the world may or may not later confirm it. This is the
subject of \S\,II. What fills a latent interstice when it is finally
confirmed is not determined by anything on the present surface, and this
paper does not claim it is; that silence is marked as a permanent open
horizon and is not crossed.

The paper proceeds as follows. \S\,II develops the present configuration
poised over an unwritten interstice and the construction of models from
the wake. \S\,III derives the re-read clock that carries distinction
across nows and states the prediction that re-reading degrades. \S\,IV
derives the decay law itself, in both its forgetting and cosmic-residue
readings, and fixes what it leaves open. \S\,V situates the paper in the
literature. The Conclusion forwards nothing beyond what is derived.

\section*{II. The Excitable Manifold and the Construction of Models}
\addcontentsline{toc}{section}{II. The Excitable Manifold and the Construction of Models}

\S\,I left the future strictly unwritten and ruled out any reach from the
present surface into the latent region. This raises a structural
question. A confirmation event requires triadic agreement among agents
before a latent interstice resolves (Paper~1, \S\,1.4); between events, the
agents bounding an unconfirmed interstice are not folding, yet they are
not inert either, since the next fold proceeds from their current
configuration. The question is what to call that state without importing
a relation to the future the substrate does not have.

The agents themselves supply the constraint. By Paper~1 \S\,1.2 the agent
manifold is immutable: agents do not change, move, or degrade. Whatever
the pre-confirmation state is, it cannot be a state of an agent, because
agents have no states to occupy. It is instead a property of the
configuration: three agents and their touch-vector signals standing in a
relation from which a fold can proceed, but in which triadic agreement
has not yet been reached. The configuration is structurally complete and
causally quiet. It does nothing until the conditions for agreement are
met, at which point it folds.

This is the structure of an excitable system: a configuration that holds
at a threshold and does not act until input crosses it, after which it
fires (van der Pol \& van der Mark, 1928; FitzHugh, 1961). We adopt the
term for the pre-confirmation configuration: the excitable manifold is
the region of agents poised to confirm a latent interstice that has not
yet folded. The term is chosen for its restriction. An excitable
configuration carries no model of what will cross its threshold and no
relation to the value the fold will produce; it carries only a release
condition. Naming the state this way records what is present --- a
threshold and the readiness to fire --- and withholds what is not: any
orientation toward the future, any anticipation, any holding-toward. The
configuration is poised, not expectant.

What looks, from outside, like a system directed at the future is a
present operation on the past. A subsystem can apply folds to its own
wake --- to structure already confirmed and sealed --- and assemble from
it a present configuration: a confirmed arrangement whose touch-vectors
point forward into the lattice (Paper~1, \S\,1.3, forward propagation
only). Such a configuration is a model in the only sense the substrate
allows: a present, fully confirmed object, built from past material, that
biases which latent interstices the now-surface is positioned to confirm
next. It selects where the manifold is excitable. It does not contain the
content of any future confirmation; it cannot, because that content is
latent and undetermined until folded. The model is past-built and
present-held, and it points at the future only in that the world may or
may not later confirm the interstices it has made the system ready for.

Two boundaries close the section. First, the construction of models is a
re-use of the wake, and re-use is a fold; by \S\,IV it sheds. A system
that builds richly from its past degrades that past in the building, by
the same law that governs all re-reading. The model-building capacity and
the decay of the material it builds from are not separate facts. Second,
what fills a latent interstice when the excitable manifold finally folds
it is not supplied by the model, by the configuration, or by anything on
the present surface. The model fixes \emph{where} the system is ready to
confirm; it does not fix \emph{what} is confirmed there. Whether the
content of the future has any determination prior to its folding is left
open, exactly as \S\,I marked it, and nothing in this section assumes an
answer.

\section*{III. The Re-read Clock}
\addcontentsline{toc}{section}{III. The Re-read Clock}

The previous sections concern a single now and the latent region ahead of
it. Time requires the relation between nows, and that relation is carried
by re-reading. A re-read is one application of \ensuremath{\hat{R}} to a
region of wake already confirmed --- the now-surface, advancing, folds a
second time on structure it or an earlier surface has written. Let
\ensuremath{n} count the folds a given region has undergone since it was
written, and let \ensuremath{f} be the rate at which those re-reads occur,
the re-reads per unit flop-time for that region. We call \ensuremath{f}
the re-read clock. It is the variable that turns the single-now quantities
of Paper~3 into across-now quantities: the form-surplus that Paper~3 sizes
at one now is shed once per fold, so \ensuremath{n} folds carry it to
\ensuremath{r^{n}} of its value, with \ensuremath{r} the per-fold survival
ratio derived in \S\,IV. The clock counts; \S\,IV gives what each count
costs.

A fold can be triggered from either side of the master equation, and the
clock has a corresponding pair of faces. A re-read induced from the bit
side --- a stressor, sensory load, the world writing over a region ---
advances the clock by overwriting; call its rate
\ensuremath{f_{\mathrm{world}}}. A re-read induced from the agent side ---
the system folding on its own wake, the model-construction of \S\,II ---
advances the same clock by re-confirmation; call its rate
\ensuremath{f_{\mathrm{self}}}. The two are faces of one operation:
\ensuremath{\hat{R}} does not record which side triggered it, and the
distinction shed per fold is the same in both cases (\S\,IV). The total
re-read rate is \ensuremath{f = f_{\mathrm{world}} + f_{\mathrm{self}}}.
The agent-side face is the same operation named, in plain register, as
recollection: to re-read a region of one's own wake is to fold it again.

One consequence follows directly and is falsifiable. Because the shed is
the same whichever face advances the clock, a region re-read from the
agent side degrades by the same law as a region overwritten from the
world side. Re-confirmation is not a free replay of a stored value; it is
a fold, and a fold sheds (\S\,IV). A memory therefore decays through use,
not only through neglect: each recollection re-reads the trace and sheds
from it, so a frequently re-read region loses fine distinction faster than
one left alone, and what is recalled is increasingly the most recent
re-reading rather than the original confirmation. This is a prediction; it
is stated here as a consequence of the shared shed and is not matched to
data in this paper.

The clock is a rate, and the present paper needs nothing more of it. Which
region a re-read falls on is set in part already and in part downstream.
When the trigger is from the bit side, the world sets the region: a
stressor or sensory load re-reads whatever it falls on. When the trigger
is from the agent side, the model of \S\,II biases the region, since a
configuration built from the wake makes some interstices excitable rather
than others. What remains --- what determines that
\ensuremath{f_{\mathrm{self}}} folds toward structure that persists across
folds rather than structure that does not --- is derived in the next
paper, and it is derived there on the shed this paper establishes:
unbound structure degrades under the per-fold shed of \S\,IV into a
configuration that cannot re-confirm, and that degradation is the
machinery the next paper's selection rule runs on. The selection is
therefore not left open but located: it is the agent-side dynamics of
Paper~5, for which \S\,IV supplies the decay. Here \ensuremath{f} enters
only as a rate.

\section*{IV. The Decay Law}
\addcontentsline{toc}{section}{IV. The Decay Law}

This section fixes what one fold costs and what many folds accumulate.
The per-fold shed is inherited: Paper~3 establishes that the confirmation
fold sheds \ensuremath{s} bits of distinction, with \ensuremath{s = 2} for
binary-agent triads, forced there and used here without re-derivation
(Paper~3, \S\,IV, Appendix~B). What this section adds is the form decay
takes across folds, and it turns on a single question the substrate
settles: whether the shed removes a fixed quantity of distinction per
fold or a fixed fraction of it.

The question is decided by what \ensuremath{\hat{R}} can act on. A re-read
operates on the region as it currently stands, not on the original
confirmation: by Axiom~4 the original is sealed in the wake, and the
re-reading surface reads the present, already-degraded copy. A fixed-budget
shed --- removing \ensuremath{s} bits of the original total each fold ---
would require the fold to act against that original total, a quantity the
surface no longer has access to. A fixed-fraction shed --- removing a set
proportion of the distinction currently present --- requires only the
present region, which is what the surface has. The fixed-fraction reading
is therefore forced, not chosen; the access constraint admits no other.
The per-fold survival ratio follows as \ensuremath{r = 2^{-s}}, and with
\ensuremath{s = 2}, \ensuremath{r = 1/4}. The derivation, with the
fixed-budget reading shown to fail, is given in Appendix~C.

Accessible distinction after \ensuremath{n} folds then decays
geometrically from its initial value toward a floor:
\begin{equation*}
D(n) = D_{\mathrm{floor}} + (D_{0} - D_{\mathrm{floor}})\, r^{n}.
\end{equation*}
The floor is not zero. By Axiom~3 the fold conserves count while shedding
distinction, so what survives unconditionally is the count-level record
that the region was confirmed --- the gist, in the agent-side register ---
while the fine distinction above the floor decays by \ensuremath{r^{n}}.
The curve is monotone decreasing in \ensuremath{n}, its form is fixed by
the access constraint rather than by a fitted rate, and it bounds below at
\ensuremath{D_{\mathrm{floor}} > 0}: detail is lost toward gist, never to
total erasure, which would require the count loss Axiom~3 forbids. The
three conditions are verified in Appendix~C. The fixed-budget reading
fails the third: a fixed subtraction drives \ensuremath{D} through the
floor to zero in finite folds, erasing the count. That the substrate
rejects fixed-budget for fixed-fraction is the same fact as the existence
of the floor.

The same per-fold shed governs a second process, at the scale of the
wake's own accumulation. A forming subsystem boundary anchors residue from
the confirmed past at a rate \ensuremath{A}, and re-reads de-anchor it: by
the law just fixed, each re-read sheds the fraction
\ensuremath{(1-r)} of the residue currently anchored, at the clock rate
\ensuremath{f}. The anchored residue \ensuremath{R} then obeys
\begin{equation*}
\frac{dR}{dt} = A - f\,(1-r)\,R,
\end{equation*}
a balance with a stable fixed point at \ensuremath{R^{*} = A / [f(1-r)]}.
The loss term carries the same \ensuremath{s = 2} as the forgetting curve,
in the same fixed-fraction sense, and reduces the residue rather than
growing it. Anchored residue therefore self-limits rather than
accumulating without bound, and the residue-to-ordinary ratio flattens as
a structural consequence rather than by tuning. The sign, the shared
\ensuremath{s}, and the commensurability of the two rates are checked in
Appendix~C.

The two processes are one law at two addresses, and the identity is sharp:
the \ensuremath{r} of the forgetting curve and the \ensuremath{r} of the
residue balance are the same \ensuremath{1/4}, used the same way ---
fixed-fraction --- for the same reason, the access constraint that forbids
a fold from reaching the original behind the present copy. Forgetting
re-reads an anchored memory; the residue balance re-reads the accumulated
wake; both are \ensuremath{\hat{R}} advancing on what is already written,
shedding the same fraction per fold. This is the unification, and it holds
structurally rather than by coincidence of magnitude. The check that the
two uses match in number and in kind is the decisive one and is carried in
Appendix~C.

What this section does not fix is also stated. The value of the
residue-to-ordinary ratio is not derived: \ensuremath{R^{*}} depends on
\ensuremath{A} and \ensuremath{f}, neither of which this paper derives, so
the flatness is structural but the magnitude is open. The identity of the
anchored residue with the material dark-matter mode, and the mass--energy
reading that would bind it to the dark-energy residual, belong to the
material side and are treated there (Volume~1, Paper~5); they are not
claimed here. And whether either curve matches measured data --- the
forgetting curve against recollection, the flattening against the observed
ratio history --- requires outside data and is not asserted. The framework
fixes the form; the world is left to confirm it.

\section*{V. Where It Sits}
\addcontentsline{toc}{section}{V. Where It Sits}

The closest metaphysical neighbor is the growing block view of time, on
which the past and present exist while the future does not, and the
present is the surface at which new moments are added to the past, so
that the block of what exists grows (Broad, 1923; Tooley, 1997; Forrest,
2004). The structure derived here is that view's structure: the wake is
the existing past, the now-surface is the surface of addition, and the
latent region is the future that does not yet exist. The fork is in
standing, not in shape. The growing block posits the asymmetry as a
metaphysical primitive; here it is derived --- the now-surface is a forced
structural object of the substrate (Paper~1), and past, present, and
future are read off the confirmed, active, and latent states rather than
asserted. The variant on which the past is sealed and only the present is
active (Forrest's dead-past version) corresponds to the now-surface being
the sole site of confirmation while the wake is locked record, with the
qualification that this paper makes no claim about where measurement is
occupied (\S\,I; the band-word discipline of Paper~3). The account departs
from eternalism, on which all times are equally real, since the future is
latent and undetermined rather than real elsewhere; and from presentism,
on which only the present exists, since the wake is real and sealed.

The closest live interlocutor on the agent side is predictive processing,
on which perception is the brain's top-down prediction of its input and
learning is the minimization of prediction error against a generative
model (Rao \& Ballard, 1999; Friston, 2010). The fork is precise and is
not a disagreement about data. On that account the present state carries a
model \emph{of} the future, projected forward; on this one the present
surface carries no content about the future at all, because the future is
latent and undetermined until folded (\S\,I), and what looks like
prediction is a present configuration built from the wake that biases
where the manifold is excitable (\S\,II). Prediction places the future
inside the present as a model; this account keeps the future unwritten and
the present merely poised. Any predictive-processing-like content, where
it exists, would live in the model-construction of \S\,II, not in the
poised state itself. The exclusion-style divergence with integrated
information theory is the subject of Paper~2 and is not reopened here.

The decay law has its nearest quantitative neighbors in two fields. In
continual learning, sequential training induces catastrophic forgetting,
and the retained-memory signal decays exponentially under ordinary
gradient descent, taking a consolidated, slower form when the weights
encoding old tasks are protected (Kirkpatrick et al., 2017); models that
combine decay with stochastic reactivation reproduce smoothly decaying
retention curves rather than a sharp cutoff. The exponential-toward-a-floor
form of \S\,IV is the same shape, and the reactivation that degrades while
it consolidates is the agent-side re-read of \S\,III seen in a network;
the correspondence is a convergence, landed after the derivation, and the
fork is that the law here is forced by one access constraint rather than
fitted as a mechanism per task. Underneath both sits the thermodynamic
floor: a logically irreversible, many-to-one operation reduces degrees of
freedom and dissipates, at least \ensuremath{kT\ln 2} per erased bit
(Landauer, 1961; Bennett, 1982). The confirmation fold is exactly such a
map --- three loci to one bit --- so the shed is the substrate's Landauer
cost, and the decay law is that cost paid once per re-read. This is noted
as the physical grounding of the shed, not as an independent result.

The poised pre-confirmation configuration corresponds to an excitable
system held below threshold (van der Pol \& van der Mark, 1928; FitzHugh,
1961), and the broader observation that information-bearing systems are
often tuned near a critical or metastable point (Bak, Tang, \&
Wiesenfeld, 1987; Beggs \& Plenz, 2003) is a convergence on the present
construction, landed after it and not used to support it: a system poised
at threshold is where its dynamic range is widest, which is the same
content the now-surface carries on the material side. The decay law also
corresponds to lossy recompression in the rate-distortion sense (Shannon,
1948): re-reading is re-encoding through a channel that cannot hold the
source's full distinction, and the gist floor is the rate below which the
code cannot compress without losing the count. The generational loss of a
copy made from a copy is the same structure observed outside the
substrate; it is noted as a correspondence, not a citation of support. The
reconstructability of the past from its forward record corresponds to the
recent result that an error-corrected history's causal memory either
propagates forward or terminates at detection (a structure independently
described for monitored many-body systems); the depth-dependent form given
here, in which rigidity grows with the wake downstream of a confirmation,
is original to this paper. On the open-future question, the account sides
structurally with the view that propositions about the unwritten future
lack a determined value until the event (the future-contingents tradition,
Aristotle, \emph{De Interpretatione} 9), but reaches it from the substrate
rather than from semantics: the value is undetermined because the latent
state is, not because of a convention about truth. The observer-relative
appearance of unpredictability --- determined from outside a boundary,
unreadable from within it --- is the substrate's account of why sensitive
dependence presents as it does to an embedded observer, and is offered as
that, not as a contribution to dynamical-systems theory.

\section*{Conclusion}
\addcontentsline{toc}{section}{Conclusion}

Time, on the substrate, is the now-surface read on both sides. The past is
the confirmed wake, sealed by causal anchoring and reconstructable from
the record it has already written, the more rigidly the deeper it lies.
The future is the latent region the surface has not folded, unwritten in
the strict sense and reachable only as a present configuration built from
the past. The present is the one surface between them, where confirmation
fires and where, alone, the record is not yet fixed. The asymmetry of past
and future is not an asymmetry of determination but of direction across
that surface: the same geometry reconstructs backward and is blocked
forward, by an undetermined draw and an unreadable exterior.

What relates one now to the next is re-reading, and re-reading sheds. The
per-fold shed of Paper~3, carried across folds by the re-read clock, fixes
a single decay law that governs the fading of a memory and the
self-limiting of the wake's residue alike, by one fraction used one way.
The law fixes the form of both and the value of neither; what is derived
is held apart from what is left open, and the open quantities are named so
that they cannot later be assumed closed.

The re-read clock has been developed here only as a rate. Its agent-side
face folds on the wake and sheds by the same law as the world's
overwriting. What determines that this face folds toward structure able to
persist across folds is taken up in the next paper, where the agent-side
dynamics are derived --- and derived on the shed established here, since it
is the per-fold degradation of \S\,IV that turns unbound structure into a
configuration that cannot re-confirm. This paper supplies the decay; the
selection that runs on it is the subject of the paper that follows. Here
the clock turns, and what it turns on decays.

\section*{Appendix A: Reconstructability of the Wake and Its Depth-Dependence}
\addcontentsline{toc}{section}{Appendix A: Reconstructability of the Wake and Its Depth-Dependence}

This appendix establishes the claim of \S\,I that a confirmed value is
reconstructable from the record its confirmation deposits downstream, that
the reconstructability increases with depth into the wake, and that it
fails at the now-surface frontier. The objects are inherited from
Volume~1; no new primitive is introduced.

\subsection*{A.1 The objects}

A confirmation event at a region resolves a latent interstice to a value
\ensuremath{\bar{B} \in \{0,1\}} with polarization
\ensuremath{\chi \equiv 2\bar{B} - 1 \in \{-1,+1\}}. By the
\ensuremath{1\to 4} multiplier (Volume~1, Paper~1, \S\,4.2), the event
emits twelve outgoing signals consumed by four downstream synthetic bits,
each carrying the source polarization \ensuremath{\chi} (Paper~1,
\S\,4.3). By Axiom~3 these signals are conserved: not created, destroyed,
or deferred. Let \ensuremath{d \geq 0} index depth into the wake, the
number of confirmed layers built downstream of the event since it was
written; \ensuremath{d = 0} is the now-surface frontier.

\subsection*{A.2 The local unit detects but does not reconstruct}

Disrupt one agent's contribution to the confirmed interstice. The agent
is immutable (Paper~1, \S\,1.2); what is disturbed is the confirmation
relation, opening the closed boundary and producing a parity syndrome
(Paper~1, \S\,5.2). The syndrome establishes that a confirmation was
present and that it was disturbed. It does not, on the local unit alone,
fix the disturbed value: a 2-simplex with one vertex contribution missing
is consistent with either \ensuremath{\chi}, since detection is of the
violation, not of the original polarization. The local unit yields
detection only.

\subsection*{A.3 The forward cone reconstructs, for \texorpdfstring{$d \geq 1$}{d >= 1}}

For \ensuremath{d \geq 1} the event has emitted its polarization into the
downstream layers, where by A.1 the signals carrying \ensuremath{\chi} are
conserved. The original \ensuremath{\chi} is therefore recoverable as the
polarization carried by the conserved downstream signals the event
itself produced, independently of the disrupted local contribution. The
number of downstream signals carrying \ensuremath{\chi} grows with the
forward cone: each layer re-emits under the same multiplier, so the count
of conserved carriers of the original polarization is non-decreasing in
\ensuremath{d}. Reconstruction is thus available for \ensuremath{d \geq 1}
and is increasingly over-determined as \ensuremath{d} grows --- the deeper
the event lies in the wake, the more independent downstream carriers fix
its value. This is the depth-dependence asserted in \S\,I.

\subsection*{A.4 The frontier limit}

At \ensuremath{d = 0} no downstream layer has yet formed, so no conserved
carrier of \ensuremath{\chi} exists outside the disrupted unit. By A.2 the
local unit yields detection only; re-confirmation must redraw the value
against the current polarization landscape, which biases the redraw
(Paper~1, \S\,4.3; Volume~1, Paper~2) without fixing it. At the frontier,
therefore, only the fact of confirmation is preserved, not its value.
Reconstructability of value is a property of the wake interior.

\subsection*{A.5 Break-test}

The result fails if reconstruction is claimed where no forward cone
carries the polarization. Two pre-registered failure conditions: (i) if
the \ensuremath{1\to 4} multiplier were content-neutral --- if downstream
signals did not inherit \ensuremath{\chi} --- then no downstream carrier
would fix the value and the claim would reduce to detection at all depths;
Paper~1 \S\,4.3 establishes the inheritance, so this does not obtain. (ii)
If the claim were extended to \ensuremath{d = 0}, it would assert
reconstruction with zero carriers, which A.4 forbids; the claim is
therefore restricted to \ensuremath{d \geq 1} and is false outside that
range, as required. The correspondence to the \ensuremath{[[4,2,2]]} code
(distance-2 detection on the local unit; reconstruction via the forward
cone) is consistent with this structure but is not used in the proof.

\section*{Appendix B: The Forward Step Is Unreadable from Within the Boundary}
\addcontentsline{toc}{section}{Appendix B: The Forward Step Is Unreadable from Within the Boundary}

This appendix establishes the second forward obstruction of \S\,I: that a
bounded subsystem cannot read the contribution its exterior makes to the
folds it will undergo, so that even a determined forward step is
unreadable from inside. It is distinct from the first obstruction, which
is the genuine undetermination of the latent draw (Axiom~0) and holds for
any observer.

\subsection*{B.1 The objects}

A subsystem is a closure: a region individuated by a boundary across which
it is defined as interior against exterior (Volume~1, Paper~3; Paper~5).
A confirmation on the boundary resolves against contributions from both
sides, since the now-surface advances on the whole local configuration,
interior and exterior alike (Paper~1, \S\,1).

\subsection*{B.2 The argument}

A measurement available to the subsystem is a confirmation it participates
in --- one whose loci lie within its closure. The exterior, by definition
of the boundary, contains loci the subsystem does not enclose and in whose
confirmations it does not participate. Their values are therefore not
among the quantities the subsystem can read. A forward fold on the
boundary resolves partly against those exterior contributions. Hence the
subsystem cannot, from within, fix the exterior input to its own next
fold, and the forward step is unreadable from inside to exactly the extent
that it depends on the exterior --- even when that step is fully determined
from a vantage that encloses both sides. The obstruction is positional: it
is a consequence of which loci fall inside the boundary, not of any
indeterminacy in the dynamics.

\subsection*{B.3 Relation to the first obstruction and to irreducibility}

The two obstructions are independent. The first removes determination at
the source (the latent draw is undetermined, Axiom~0); the second removes
readability from a position (the exterior is determined but unenclosed).
Either alone blocks forward reading; together they fix that prediction
from within runs against both an undetermined draw and an unreadable
exterior. The positional obstruction is the time-like reading of the
record's computational irreducibility (Volume~1, Paper~1, \S\,7.2a): the
record cannot be compressed, and the forward step cannot be shortcut from
within.

\subsection*{B.4 Break-test}

The claim fails if an interior measurement could fix an exterior locus. By
B.1 the boundary is defined precisely by non-enclosure of the exterior, so
an interior confirmation enclosing an exterior locus would contradict the
definition of the closure; no such measurement exists for a genuine
boundary. The claim is correspondingly vacuous for a system with no
exterior, which is the total-system case treated separately (Paper on the
total system); there the obstruction does not arise because there is no
unenclosed remainder, consistent with the result being positional.

\section*{Appendix C: The Decay Law}
\addcontentsline{toc}{section}{Appendix C: The Decay Law}

This appendix derives the per-fold survival ratio, the forgetting curve
and its three conditions, the cosmic-residue balance and its sign and
fixed point, and the unification check that the two uses share one ratio
used one way. The per-fold shed \ensuremath{s} is inherited (Paper~3,
Appendix~B), with \ensuremath{s = 2} for binary-agent triads.

\subsection*{C.1 Fixed fraction is forced; the survival ratio}

A re-read applies \ensuremath{\hat{R}} to a region as it currently stands.
By Axiom~4 the original confirmation is sealed and cannot be re-accessed;
the re-reading surface acts on the present copy. A fixed-budget shed
removes \ensuremath{s} bits of the original total per fold and so requires
the original total as an operand, which the surface does not possess. A
fixed-fraction shed removes a set proportion of the distinction presently
accessible and requires only the present region. Only the latter is a
substrate-available operation; the fixed-fraction reading is forced. The
per-fold survival ratio is \ensuremath{r = 2^{-s}}, and for \ensuremath{s
= 2}, \ensuremath{r = 1/4}. The value tracks \ensuremath{s}: it is fixed
by the inherited shed, not chosen to resemble decay.

\subsection*{C.2 The forgetting curve and its three conditions}

With a fixed fraction \ensuremath{r} surviving per fold, accessible
distinction after \ensuremath{n} folds is
\ensuremath{D(n) = D_{\mathrm{floor}} + (D_{0} - D_{\mathrm{floor}}) r^{n}}.
(A1) Monotone: \ensuremath{r \in (0,1)} gives \ensuremath{r^{n}} strictly
decreasing, so \ensuremath{D(n)} decreases monotonically toward
\ensuremath{D_{\mathrm{floor}}}; shed distinction is not regained, since
Axiom~4 seals the original. (A2) Derived form: the exponential form
follows from the fixed-fraction result of C.1, and the rate
\ensuremath{r = 2^{-s}} from the inherited \ensuremath{s}; no rate is
fitted. (A3) Floor: by Axiom~3 the fold conserves count, so the count-level
record (gist) is never shed; \ensuremath{D_{\mathrm{floor}} > 0} and the
curve distinguishes detail-loss from total erasure. The fixed-budget
reading fails A3: a fixed subtraction reaches and crosses zero in finite
folds, destroying the count, which Axiom~3 forbids. The floor existing and
fixed-fraction being forced are the same fact.

\subsection*{C.3 The cosmic-residue balance}

A forming boundary anchors residue at rate \ensuremath{A} and de-anchors
it by re-reading at clock rate \ensuremath{f}, each re-read shedding the
fraction \ensuremath{(1-r)} of the residue then anchored. Hence
\ensuremath{dR/dt = A - f(1-r)R}. (B1) Sign: for
\ensuremath{R, f, (1-r) > 0} the loss term is negative, reducing the
residue, opposing accumulation. (B2) Same shed: \ensuremath{(1-r) = 1 -
2^{-s}} with the same \ensuremath{s = 2}, giving \ensuremath{3/4}. (B3)
Units: \ensuremath{A} is residue per time; \ensuremath{f(1-r)R} is (folds
per time)(dimensionless)(residue) \ensuremath{=} residue per time;
commensurable, so cancellation is well-posed. The fixed point
\ensuremath{R^{*} = A/[f(1-r)]} is stable, since
\ensuremath{\partial_{R}[A - f(1-r)R] = -f(1-r) < 0}; the residue
self-limits and the ratio flattens structurally.

\subsection*{C.4 The unification check}

The decisive condition (X1): the \ensuremath{r} of C.2 and the
\ensuremath{r} of C.3 must be the same value used the same way. Both are
\ensuremath{r = 2^{-s}} with \ensuremath{s = 2}; both enter as fixed
fractions, not fixed budgets; and both are fixed-fraction for the same
reason --- the access constraint of C.1, which forbids a fold from
reaching the original behind the present copy, in the memory region and in
the accumulated wake alike. The unification is structural, not a numerical
coincidence between two independently tuned processes.

\subsection*{C.5 What is not fixed; break-test}

The value of \ensuremath{R^{*}} depends on \ensuremath{A} and
\ensuremath{f}, neither derived here; the flatness is structural and the
magnitude is open. The break-test is X1 itself: had forgetting forced
fixed-fraction while the residue balance forced fixed-budget, or had the
two required different \ensuremath{s}, the single-law claim would be false
even with each half internally sound. Both forced fixed-fraction from one
access constraint and share \ensuremath{s = 2}; the test is passed. A
second pre-registered failure: if accessible distinction were defined
against the sealed original rather than the present copy, C.1 would admit
fixed-budget and A3 would fail; Axiom~4 forbids that access, closing the
route.

\section*{Acknowledgments}
\addcontentsline{toc}{section}{Acknowledgments}

This paper develops the now-surface of Paper~1 and the per-fold shed of
Paper~3, and inherits its primitives from Volume~1. It is one paper of a
larger derivation campaign and does not stand alone.

\section*{References}
\addcontentsline{toc}{section}{References}

\begin{reflist}
Bak, P., Tang, C., \& Wiesenfeld, K. (1987). Self-organized criticality: An explanation of the 1/f noise. \emph{Physical Review Letters}, 59(4), 381--384. doi:10.1103/PhysRevLett.59.381.

Beggs, J. M., \& Plenz, D. (2003). Neuronal avalanches in neocortical circuits. \emph{The Journal of Neuroscience}, 23(35), 11167--11177. doi:10.1523/JNEUROSCI.23-35-11167.2003.

Bennett, C. H. (1982). The thermodynamics of computation---a review. \emph{International Journal of Theoretical Physics}, 21(12), 905--940. doi:10.1007/BF02084158.

Broad, C. D. (1923). \emph{Scientific Thought}. London: Kegan Paul, Trench, Trubner \& Co.

FitzHugh, R. (1961). Impulses and physiological states in theoretical models of nerve membrane. \emph{Biophysical Journal}, 1(6), 445--466. doi:10.1016/S0006-3495(61)86902-6.

Forrest, P. (2004). The real but dead past: A reply to Braddon-Mitchell. \emph{Analysis}, 64(4), 358--362. doi:10.1093/analys/64.4.358.

Friston, K. (2010). The free-energy principle: A unified brain theory? \emph{Nature Reviews Neuroscience}, 11(2), 127--138. doi:10.1038/nrn2787.

Kirkpatrick, J., Pascanu, R., Rabinowitz, N., Veness, J., Desjardins, G., Rusu, A. A., Milan, K., Quan, J., Ramalho, T., Grabska-Barwinska, A., Hassabis, D., Clopath, C., Kumaran, D., \& Hadsell, R. (2017). Overcoming catastrophic forgetting in neural networks. \emph{Proceedings of the National Academy of Sciences}, 114(13), 3521--3526. doi:10.1073/pnas.1611835114.

Landauer, R. (1961). Irreversibility and heat generation in the computing process. \emph{IBM Journal of Research and Development}, 5(3), 183--191. doi:10.1147/rd.53.0183.

Rao, R. P. N., \& Ballard, D. H. (1999). Predictive coding in the visual cortex: A functional interpretation of some extra-classical receptive-field effects. \emph{Nature Neuroscience}, 2(1), 79--87. doi:10.1038/4580.

Tooley, M. (1997). \emph{Time, Tense, and Causation}. Oxford: Oxford University Press.

van der Pol, B., \& van der Mark, J. (1928). The heartbeat considered as a relaxation oscillation, and an electrical model of the heart. \emph{The London, Edinburgh, and Dublin Philosophical Magazine and Journal of Science}, 6(38), 763--775. doi:10.1080/14786441108564652.
\end{reflist}


\end{document}
